% Pursuing Stacks (cotes 134-2 to 134-8): archive cote 134-6, site page 7 (PDF page 8), typescript page 326.
% typescript. Draft compared: ps5.tex:1682-1743.
% AI-transcribed by Claude Opus 5.5 (claude-opus-5-5, adaptive thinking, effort high), Message Batches API, batch ps3 (msgbatch_01XdFmuVVBAcgc7EFxdtCqsj), 2026-11-12; not yet checked by eye.
\scanpage{134-6/7}
\mspage{326}

an\xout{xxxx}\inserted[typed]{aspheric} functor (10), in such a way \xout{moreover} that it factors the canoni-%
cal functor (6) (granting Whitehead's theorem holds for A), and more-%
over that for an aspheric functor (10') j: B \hand{$\longrightarrow$} A\^{} , giv\xout{x}ing rise
to
\[ \overline{j^{*}_{ab}} :\ \mathrm{Hotab}_{A} \hand{$\longrightarrow$} \mathrm{Hotab}_{B} \quad , \]
the corresponding diagram
\diagram{
\begin{tikzcd}[column sep=small]
\mathrm{Hotab}_{A} \arrow[rr, "\hand{}"] \arrow[dr, "\hand{}"'] & & \mathrm{Hotab}_{B} \arrow[dl, "\hand{}"] \\
& \mathrm{Hotab} &
\end{tikzcd}
}
(15)

should commute, where the vertical arrows are the functors (14).

For defining (14) in this general case, recalling that
\[ \mathrm{Hotab} \hand{$\xleftarrow{\approx}$} \mathrm{Hotab}_{\hand{$\Delta$}} \quad , \]
we can't help it and have to use diagram (3) and the functor
\[ i_{A} : A^{\wedge} \hand{$\longrightarrow$} (\mathrm{Cat}) \quad , \]
or rather ((Cat) serving only as an intermediary) the functor
\[ u \overset{\mathrm{def}}{=} i^{*} i_{A} : A^{\wedge} \hand{$\longrightarrow$} \hand{$\Delta$}^{\wedge} \quad , \]
so that (6) can be viewed as \xout{\unsure{the composition}} deduced \inserted{by localization of $A^{\wedge}$} from the
composition
\[ (\hand{$*$}) \qquad A^{\wedge} \hand{$\xrightarrow{u}$} \hand{$\Delta$}^{\wedge} \hand{$\xrightarrow{Wh_{\Delta}}$} \hand{$\Delta$}^{\wedge}_{ab} \hand{$\longrightarrow$} \mathrm{Hotab}_{\hand{$\Delta$}} \quad ( \hand{$\xrightarrow{\approx}$} \mathrm{Hotab} ) \quad . \]
\xout{The}\inserted[typed]{One} difficulty here is that u does not commute to finite prod\xout{d}ucts,
and hence does not induce a functor\xout{,} \xout{it would seem, from $A^{\wedge}_{ab}$}
\[ A^{\wedge}_{ab} \hand{$\longrightarrow$} \hand{$\Delta$}^{\wedge}_{ab} \quad , \]
it would seem. Now this difficulty, I just noticed, can be overcome,
using the fact that i$_{A}$ and hence also u \uline{commutes to fibered products},
or, what amount to the same, induces an \uline{exact} functor
\[ u_{o} :\ A^{\wedge} \hand{$\longrightarrow$} \hand{$\Delta$}^{\wedge}_{/F} \ \deleted{\illegible{1}}\inserted{$\approx$} (\hand{$\Delta$}_{/F})^{\wedge} \hand{,} \]
where
\[ F = u(e_{A^{\wedge}}) = i^{*}(A) \]
is a suitable object in \hand{$\Delta$}\^{}. A fortiori u$_{o}$ commutes to finite pro-%
ducts, hence transforms abelian group objects into same i.e. induces
\[ u_{oab} :\ A^{\wedge}_{ab} \hand{$\longrightarrow$} (\hand{$\Delta$}_{/F})^{\wedge}_{ab} \quad , \]

\margin{7}
